\section{Considerações Finais}

A implementação e subsequente estresse da arquitetura analítica do CCO-DataOps consolidam a eficácia das modernas engrenagens de \textit{DataOps} na gestão de crises ambientais e infraestrutura de missão crítica. Através do acoplamento do Apache Superset ao motor OLAP do Apache Druid, demonstrou-se, com respaldo em metodologias laboratoriais, a viabilidade técnica de extrair \textit{insights} imediatos do cruzamento de bilhões de dados telemétricos geoespaciais com dicionários complexos do território brasileiro.

O arcabouço tecnológico provou ser resiliente e perfeitamente contido no espectro de custos de computação em nuvem moderada. Evidenciou-se que a adoção precoce de compactação profunda (\textit{Storage Footprint} otimizado) e varreduras colunares anula a necessidade de dispêndios orçamentários volumosos em armazenamento. Além disso, as descobertas obtidas pelo isolamento das anomalias de \textit{performance} fornecem aos Engenheiros de Dados do projeto diretrizes clarividentes --- como a imperiosa migração de \textit{Joins} onerosos para etapas puras de \textit{ETL} via Apache Spark.

Em última instância, este relatório atesta não somente a capacidade atual da plataforma de operar em cenários limítrofes, mas, fundamentalmente, firma o alicerce científico e arquitetural necessário para escalar o monitoramento de Focos de Incêndio a níveis governamentais absolutos, unindo dados abertos de saúde, logística e segurança à fronteira da visualização analítica geoespacial.
